%% ma: L12-B08-0106
%% lop: 12
%% chuong: 2
%% bai: 8
%% dang: 12.8.27
%% loai: TN4
%% muc: TH
%% dap_an: "A"
%% nguon: "CĐ2.3.pdf giáo viên gửi 27/10/2026 (M07, Lê Minh Khôi – Nguyễn Anh Khôi, Chương II Chủ đề 3 Biểu thức tọa độ của các phép toán vectơ), Câu 17"
%% kien_thuc: "Bài 8 (tọa độ trọng tâm tam giác), Bài 6 (tính chất trọng tâm)"
%% trang_thai: chinh-thuc
\begin{ex}
Trong không gian $Oxyz$, cho ba điểm $A(16;11;8)$, $B(2;12;7)$, $C(0;1;9)$ và điểm $M$ bất kỳ. Tọa độ của điểm $P$ thỏa mãn $\dfrac13\cdot\left(\overrightarrow{MA}+\overrightarrow{MB}+\overrightarrow{MC}\right)=\overrightarrow{MP}$ là
\choice{\True $P(6;8;8)$}{$P(9;12;12)$}{$P(3;4;3)$}{$P(4;3;3)$}
\loigiai{Đẳng thức đúng với mọi $M$ nên $P$ là trọng tâm tam giác $ABC$: $P\left(\dfrac{18}3;\dfrac{24}3;\dfrac{24}3\right)=(6;8;8)$.}
\end{ex}
